\chapter*{Notation and Conventions}
\addcontentsline{toc}{chapter}{Notation and Conventions}
\begin{tabularx}{\textwidth}{@{}L{0.28\textwidth}Y@{}}
\toprule
Convention & Meaning \\
\midrule
\(\jj\) & Imaginary unit. Electrical engineering uses \(\jj\) because \(i\) usually denotes current. \\
\(v(t), i(t)\) & Time-domain voltage and current. \\
\(V, I\) & Phasor voltage and current. Phasor magnitudes are RMS unless stated otherwise. \\
\(\omega=2\pi f\) & Angular frequency in radians per second is distinct from frequency in hertz. \\
\(Z=R+\jj X\) & Impedance in rectangular form. Positive \(X\) is inductive; negative \(X\) is capacitive. \\
Passive sign convention & Current enters the terminal marked positive for absorbed power. \\
\bottomrule
\end{tabularx}

Capacitor voltage and inductor current are continuous under finite,
non-impulsive excitation. Transfer functions normally assume zero initial
conditions. Phasor analysis applies to sinusoidal steady state.
